-------------------------------------------------------------------
\chapter*{Abstract}
%-------------------------------------------------------------------

In industrial automation, image recognition plays an important role. It is widely used in industries in simple forms for contour detection and as the usage is booming, there are a lot of use cases available.

In the context of modern manufacturing and recycling industries, efficient and accurate identification of patterns in cardboard materials is critical for quality control, sorting, and material reuse. This project explores the application of machine learning and \ac{AI} techniques to develop a robust pattern recognition system tailored specifically for cardboard. Leveraging \ac{CNNs} and advanced image processing algorithms, the system is designed to detect, classify, and analyze various patterns that commonly occur in cardboard materials, including texture variations, not distinguishable even by human eye.

The project investigates different machine learning models, evaluates their performance, and optimizes them for accuracy and computational efficiency. A comprehensive dataset of cardboard images is curated and annotated to train and validate the models. The outcomes of this project demonstrate the potential of AI-driven solutions to significantly enhance the efficiency and reliability of pattern recognition tasks in the cardboard industry, paving the way for more sustainable and automated manufacturing processes.


\thispagestyle{empty}
\mbox{}
\newpage

